\section{Assembly of the lifted systems}
\label{sec:lifted_system}

\subsection{First closure}

The first lifted state is
\[
z_2^{\mathcalE}(x)
=
\col(x,\Phi_2^{\mathcalE}(x))
\in\R^{26}.
\]
The original state equations are exactly linear in $x$ and the infection block
of $\Phi_2^{\mathcalE}$. Differentiating the quadratic observables gives terms
retained in $\Phi_2^{\mathcalE}$ plus cubic terms. Thus
\[
\dot z_2^{\mathcalE}
=
\matK_2^{\mathcalE}z_2^{\mathcalE}
+r_2^{\mathcalE}(x),
\]
where the first 12 residual components are zero. For a channel
$(i,j)\in\mathcalE$, the two nonzero residual blocks have components
\[
(r_{SI}^{\mathcalE})_{(i,j)}
=
-\beta\sum_{\ell\in\Min(i)}
 a_{i\ell}S_iI_\ell I_j,
\]
\[
(r_{SE}^{\mathcalE})_{(i,j)}
=
-\beta\sum_{\ell\in\Min(i)}
 a_{i\ell}S_iI_\ell E_j
+
\beta\sum_{\ell\in\Min(j)}
 a_{j\ell}S_iS_jI_\ell.
\]
The 42 cubic seeds identified in \cref{sec:cubic_generation} are precisely the
distinct monomials appearing in these residual components.

\subsection{Order-three state}

Choose the deterministic ordering in \cref{app:dictionary} and define
\[
\Phi_3^{\mathcalE}(x)
=
[q(x)]_{q\in\mathcal M_3^{\mathcalE}}
\in\R^{67}.
\]
The order-three lifted state is
\[
z_3^{\mathcalE}(x)
=
\col(x,\Phi_2^{\mathcalE}(x),\Phi_3^{\mathcalE}(x))
\in\R^{93}.
\]
Its dimension is
\[
\dim z_3^{\mathcalE}
=12+14+67
=93.
\]

The exact infinite hierarchy has the degree-block structure
\[
\begin{aligned}
\dot x
&=\matK_{1,1}x+\matK_{1,2}\Phi_2^{\mathcalE},\\
\dot\Phi_2^{\mathcalE}
&=\matK_{2,2}\Phi_2^{\mathcalE}
 +\matK_{2,3}\Phi_3^{\mathcalE},\\
\dot\Phi_3^{\mathcalE}
&=\matK_{3,3}\Phi_3^{\mathcalE}
 +\matK_{3,4}\Phi_4^{\mathcalE}.
\end{aligned}
\]
The relevant retained block sizes are
\[
\matK_{1,1}\in\R^{12\times12},
\quad
\matK_{1,2}\in\R^{12\times14},
\]
\[
\matK_{2,2}\in\R^{14\times14},
\quad
\matK_{2,3}\in\R^{14\times67},
\quad
\matK_{3,3}\in\R^{67\times67}.
\]

At order three, the quartic block is not retained. Therefore
\[
\dot z_3^{\mathcalE}
=
\matK_3^{\mathcalE}z_3^{\mathcalE}
+r_3^{\mathcalE}(x),
\]
with
\[
r_3^{\mathcalE}(x)
=
\col\bigl(0_{12},0_{14},\LieQ\Phi_3^{\mathcalE}(x)\bigr).
\]
Every nonzero residual component is homogeneous of degree four. No omitted
cubic term remains because $\mathcal M_3^{\mathcalE}$ is linearly closed.

\subsection{Comparison with complete degree-three liftings}

A complete ordered lifting through degree three would contain
\[
12+12^2+12^3=1884
\]
coordinates. A complete symmetric monomial lifting through degree three would
contain
\[
12+\binom{13}{2}+\binom{14}{3}
=12+78+364
=454
\]
coordinates. The graph-induced linearly closed construction contains 93.
This comparison is finite-dimensional and illustrative; it does not assert
minimality or optimality of the graph-induced dictionary.
